\chapter{Method}
\label{chap:method}

\section{Objective}

The goal of this work is to detect, on an edge device, when an obstacle blocks
an evacuation route (a door or a corridor). Among the pipelines compared in this
thesis, this chapter develops the third one,
\emph{Unsupervised Anomaly Detection with PatchCore} (P3), whose distinctive
property is that it learns only from normal, unobstructed scenes and requires no
labelled examples of obstructions. The following sections derive the components
of the method, their mathematical formulation, and the data on which they
operate; implementation choices and numerical settings are deferred to
Chapter~\ref{chap:setup}, and the empirical evaluation to
Chapter~\ref{chap:analysis}.

\section{Pipeline Overview}

\subsection{Unsupervised Anomaly Detection with PatchCore (P3)}

\subsubsection{Motivation and design rationale}

The detection of emergency-exit occlusions presents a fundamental labelling
bottleneck. Normal scenes - clear, unobstructed doorways and corridors - are
easy to capture at scale, but the space of possible obstructions is open-ended:
wheelchairs, carts, stretchers, waste bins, cardboard boxes, and any combination
thereof can occupy a doorway in arbitrary configurations, orientations, and
lighting conditions. Collecting and labelling a representative sample of
anomalous scenes is time-consuming, and any finite labelled set risks leaving
coverage gaps on obstacle types not seen during training.

P3 sidesteps this bottleneck by framing occlusion detection as
\emph{one-class anomaly detection}: the system learns only from normal scenes
and, at inference time, flags any image whose feature representation deviates
significantly from what was memorised as normal. This removes the need for
labelled anomalous examples entirely. The choice to adapt PatchCore
\cite{Roth2022}, originally proposed for industrial surface inspection, to the
emergency-exit domain is motivated by three properties that align with the
deployment context:

\begin{enumerate}
  \item \textbf{One-shot operation}: a single reference photograph of each
  monitored location is sufficient to build the entire anomaly model. In a real
  deployment, each fixed camera captures one reference frame of the clear exit;
  no further data collection is required.
  \item \textbf{Spatial resolution}: PatchCore operates at the patch level rather
  than producing a single image-level embedding, preserving the spatial
  structure of the anomaly and enabling heatmap visualisation and
  spatially-aware filtering.
  \item \textbf{Edge compatibility}: inference - feature extraction through a
  frozen backbone followed by a nearest-neighbour lookup in a compact memory -
  runs entirely on CPU, requires no gradient computation, and can be
  parallelised per camera without shared state.
\end{enumerate}

Among the four pipelines, P3 is the most data-efficient: it needs only a single
normal reference image and a pre-trained backbone, with no annotation effort to
deploy on a new camera. The trade-off is that it cannot identify \emph{which}
obstacle is present - it only signals that something is anomalous - and it is, in
principle, sensitive to any scene change, including non-obstacle changes such as
lighting variations, shadows, or background movement. Managing these
false-positive sources is the central methodological challenge of this chapter.

\subsubsection{Key intuition}

The core intuition is that the intermediate representations of a deep
convolutional network pre-trained on a large natural-image corpus are broadly
transferable: they capture generic visual patterns (textures, edges, object
parts) that remain discriminative across domains, including surveillance footage
of building interiors. A patch belonging to a normal scene of a fire-exit door
yields a feature vector close to other normal patches in this embedding space,
whereas a patch belonging to an unexpected obstacle - the wheel of a cart, the
fabric of a stretcher, the surface of a box - yields a feature vector distant
from all stored normal representations, hence a high anomaly score.

The memory is built from augmented variants of a single reference image rather
than from a large dataset, exploiting the fixed-camera assumption: the camera
does not move, so the background is stable across days and the only source of
variation is illumination. Photometric augmentation at build time teaches the
memory to tolerate this variation, making the decision threshold robust to
normal fluctuations without requiring multiple observation days.

It is worth stressing why the alternative - populating the memory with a
collection of \emph{similar} exit scenes (many different corridors, many
different doors) - would be both infeasible and counterproductive. Such scenes are
visually too heterogeneous: a memory mixing dozens of unrelated corridors and
doorways would model the generic appearance of ``a corridor'' or ``a door'',
which is precisely the wrong target. Its feature manifold would be so broad that
an obstacle in any specific scene could fall close to some unrelated normal scene
in the bank and be scored as normal, destroying the detector's sensitivity. The
task is not to recognise corridors and doors in general, but to recognise
\emph{this} particular exit as it normally looks. This is the natural consequence
of the deployment model, in which each camera monitors exactly one exit route:
the normal class for that camera is its own fixed scene, and the only relevant
variability is the illumination and the background change of that scene over time. A per-camera memory
built from one reference image is therefore not a limitation imposed by scarce
data, but the correct modelling choice for the objective.

\subsubsection{Feature extraction}

Let $I_{\mathrm{ref}}$ denote the normal reference image of a monitored location.
Every image is resized to a fixed resolution and passed through a frozen
convolutional backbone $f$. Following the original PatchCore formulation,
features are read from two intermediate blocks, $\ell_2$ and $\ell_3$: the
shallower block preserves local detail while the deeper block carries semantic
context, and their combination balances the two. The shallowest block encodes
mostly low-level texture with little semantic content, and the deepest block is
too abstract and biased toward the pre-training classes; both are therefore
excluded.

The deeper feature map is upsampled bilinearly to the spatial resolution of the
shallower one and concatenated along the channel axis, yielding a joint tensor of
patch vectors arranged on a $h \times w$ grid. A local average pooling over a
$3 \times 3$ neighbourhood (the neighbourhood aggregation of
\cite{Roth2022}) is then applied: each patch vector is replaced by the mean of
its $3 \times 3$ spatial neighbourhood,

\begin{equation}
  \phi(x) \;=\; \frac{1}{|\mathcal{N}_3(x)|} \sum_{x' \in \mathcal{N}_3(x)} \tilde\phi(x'),
  \label{eq:nbagg}
\end{equation}

where $\tilde\phi(x')$ is the raw concatenated feature at grid location $x'$ and
$\mathcal{N}_3(x)$ is the $3 \times 3$ window centred at $x$. This enlarges the
receptive field of each vector and reduces sensitivity to pixel-level noise. The
result is a set of $N_p = h \cdot w$ patch feature vectors per image.

\subsubsection{Build phase}

The build phase constructs the \emph{memory bank} $\mathcal{M}$, the finite set of
normal patch features against which every test patch will later be compared. It
proceeds in three steps - augmentation, feature accumulation, and coreset
reduction - each of which is detailed below. The guiding principle is that
$\mathcal{M}$ must represent the full appearance of the normal scene under its
expected illumination variation, yet remain small enough for a real-time
nearest-neighbour search on an edge device.

The build phase is a \emph{one-off, offline} operation, performed once when the
camera is installed and its reference frame $I_{\mathrm{ref}}$ is captured. The
resulting memory bank is then stored and reused unchanged for every subsequent
inference; nothing in the build phase runs at monitoring time. A rebuild is
required only when the \emph{normal} reference scene itself changes in a way that
is not a transient occlusion - for example a permanent rearrangement of the
furniture in view, a repainting, or a change of the camera's mounting - so that
the old $I_{\mathrm{ref}}$ no longer describes the normal state. Routine
illumination changes and background variation over time are deliberately
\emph{not} a reason to rebuild: they are exactly what the augmentation and the
calibration are designed to absorb. This separation between a rare, offline build
and a continuous, lightweight inference is what makes the method practical on an
edge device.

\paragraph{Step 1 - Augmentation.} A single reference frame is too
narrow a description of the normal scene: it captures one instant of illumination,
whereas the same fixed camera will, over days, see the scene under brighter or
dimmer light, warmer or cooler colour temperature, and varying sensor noise. To
let the memory tolerate this variation, $n_{\mathrm{aug}}$ augmented variants are
generated from $I_{\mathrm{ref}}$, where $n_{\mathrm{aug}}$ is the number of bank
variants. The augmentations are predominantly photometric - brightness, contrast,
saturation, colour-temperature shifts, additive noise, and a mild blur - each
motivated by a physical source of normal variation (time of day, light type,
sensor noise, slight defocus). \emph{Large} geometric transforms (flips, wide
rotations, crops) are deliberately avoided: the camera is physically fixed, so
orientation is semantically meaningful and must not be treated as a source of
invariance, otherwise the memory would also accept an obstacle appearing in a
flipped or arbitrarily relocated position as normal.

A controlled exception is a \emph{small} geometric perturbation, off by default
and enabled only for cameras subject to micro-vibration. When active, each
variant is additionally shifted by a random translation of at most
$\delta_{\max}$ pixels on each axis and rotated by a random angle of at most
$\rho_{\max}$ degrees, with $\delta_{\max}$ of a few pixels and $\rho_{\max}$ of a
couple of degrees. Its purpose is opposite to that of the forbidden large
transforms: rather than making the memory invariant to where an obstacle sits, it
makes the memory tolerant to a slight global displacement of the \emph{whole}
fixed scene, as produced by mechanical vibration or a small mounting drift. Such a
displacement moves every patch by the same tiny amount, so a few-pixel jitter
keeps a normal scene normal without ever spanning the distance an obstacle would
cover. This is the cheapest of the available robustness measures, since it acts
entirely at build time and adds no inference-time cost. 

\paragraph{Step 2 - Feature accumulation.} Each of the $n_{\mathrm{aug}}$ variants
is passed through the feature pipeline of equation~\eqref{eq:nbagg}, producing
$N_p$ patch vectors per variant. Stacking the patches of all variants gives the
\emph{raw memory}

\begin{equation}
  \mathcal{M}_{\mathrm{raw}} \;=\;
  \bigcup_{v=1}^{n_{\mathrm{aug}}} \big\{\, \phi^{(v)}(x) \;:\; x \in \text{grid} \,\big\},
  \qquad
  |\mathcal{M}_{\mathrm{raw}}| = n_{\mathrm{aug}} \cdot N_p ,
  \label{eq:rawmemory}
\end{equation}

where $\phi^{(v)}(x)$ is the aggregated feature at grid location $x$ of variant
$v$. This mirrors the memory construction of \cite{Roth2022},
$\mathcal{M} = \bigcup_{x_i \in \mathcal{X}_N} \mathcal{P}(\phi(x_i))$, with the
nominal training set $\mathcal{X}_N$ here instantiated by the augmented variants
of the single reference $I_{\mathrm{ref}}$. Each spatial location of the scene
thus contributes $n_{\mathrm{aug}}$ nearby vectors - one per illumination variant -
which together trace out the small region of feature space that a normal patch at
that location can occupy. This is exactly the tolerance the detector needs: at
test time a normal patch, under any illumination within the augmented range, will
fall close to at least one of these vectors and receive a low anomaly score.

\paragraph{Step 3 - Coreset reduction.} The raw memory is highly redundant:
adjacent grid locations on flat surfaces (a wall, a floor, a door panel) and the
$n_{\mathrm{aug}}$ illumination copies of each location produce many near-duplicate
vectors. Storing all of them would inflate both the memory footprint and the
nearest-neighbour cost without improving coverage. The memory is therefore
compressed to a coreset $\mathcal{M} \subset \mathcal{M}_{\mathrm{raw}}$ of size a
fraction $p$ of the original - the \emph{coreset ratio} - that minimises the worst
case distance of any raw vector to its closest retained vector
\cite{SenerSavarese2018,Roth2022}:

\begin{equation}
  \mathcal{M}^{\ast} \;=\;
  \arg\min_{\mathcal{M} \subset \mathcal{M}_{\mathrm{raw}}}
  \; \max_{m \in \mathcal{M}_{\mathrm{raw}}}
  \; \min_{m' \in \mathcal{M}} \lVert m - m' \rVert_2 .
  \label{eq:coreset-obj}
\end{equation}

This minimax (facility-location) objective is NP-hard, so it is approximated by a
greedy farthest-point procedure: starting from one random point, the algorithm
repeatedly adds the still unselected vector that is \emph{farthest} from the
current selection,

\begin{equation}
  m_{t} \;=\; \arg\max_{m \in \mathcal{M}_{\mathrm{raw}} \setminus \mathcal{M}_t}
  \; \min_{m' \in \mathcal{M}_t} \lVert m - m' \rVert_2 ,
  \qquad
  \mathcal{M}_{t+1} = \mathcal{M}_t \cup \{m_t\} .
  \label{eq:coreset}
\end{equation}

Choosing, at each step, the point that maximises its minimum distance to the
already-selected set spreads the retained vectors as evenly as possible over the
occupied region of feature space: dense clusters of near-duplicates are summarised
by a few representatives, while isolated points - which carry unique information
about the normal scene - are kept. The procedure stops when
$|\mathcal{M}| = \lceil p \cdot |\mathcal{M}_{\mathrm{raw}}| \rceil$, subject to a
small floor on the number of retained points so that even a very low $p$ leaves a
usable bank. Because evaluating equation~\eqref{eq:coreset} directly in the full
feature dimension is expensive, the distance comparisons are performed on a
low-dimensional random projection $\psi$ that approximately preserves pairwise
Euclidean distances (a Johnson-Lindenstrauss embedding \cite{enwiki:1357719950}), exactly as in
\cite{Roth2022}, while the selected vectors are stored at their original
dimension. The result is a compact bank that
fits in a few megabytes and enables sub-millisecond nearest-neighbour lookups, a
prerequisite for edge deployment, while retaining near-uniform coverage of the
normal-feature manifold.

\subsubsection{Threshold calibration}

\paragraph{Design space.} Several threshold strategies are possible for a
one-class detector: a fixed global value, a per-camera percentile, a
ROC-optimal cut calibrated on labelled anomalies, or a parametric rule derived
from the distribution of normal scores. A fixed global threshold requires that
the absolute scale of anomaly scores be consistent across cameras, which does
not hold in practice - a plain opaque door scores structurally lower than a glass
door with a visible outdoor background, so a threshold tuned on one would
misfire on the other. A ROC-optimal threshold requires labelled anomalies,
violating the one-class premise and risking overfitting to the appearance of the
synthetic anomalies used for tuning. A per-camera percentile is more sensitive to
outliers in the small calibration sample. The chosen strategy is a parametric
rule that adapts to each scene using only normal data.

\paragraph{Parametric rule.} For each of $n_{\mathrm{cal}}$ additional augmented
variants of $I_{\mathrm{ref}}$ (the \emph{calibration variants}, disjoint from the
bank variants and generated with an independent random seed), the maximum patch
anomaly score over the image is computed. These $n_{\mathrm{cal}}$ scores form the
empirical distribution of the maximum anomaly score of a normal scene, with mean
$\mu$ and standard deviation $\sigma$. The decision threshold is

\begin{equation}
  \theta \;=\; \mu + k\,\sigma ,
  \label{eq:threshold}
\end{equation}

where $k$ is a sensitivity hyperparameter. This is the rule of statistical
process control (Shewhart charts): an event is flagged when it lies $k$ standard
deviations above the normal mean. Under a Gaussian assumption, $k=3$ corresponds
to a theoretical false-positive rate of about $0.13\%$ and $k=4$ to about
$0.003\%$; the calibration scores are not exactly Gaussian (they are maxima of
correlated patch scores), but $k$ still provides a monotone, interpretable
control over the false-positive/recall trade-off. It is the sole
deployment-time hyperparameter and encodes the asymmetric cost of a missed
obstruction relative to a false alarm: for safety-critical exits $k$ is set low
($2.5$ to $3.0$) to favour recall, and higher ($4.0$ to $5.0$) where false
alarms are operationally costly.

The original $I_{\mathrm{ref}}$ is excluded from calibration as well: being the
closest possible image to the memory (which is built from its augmentations),
its score is structurally the minimum of the distribution rather than a
representative sample, and including it would bias $\mu$ downward. Reserving
$I_{\mathrm{ref}}$ as the actual normal test point - unseen by both the memory and
the calibration - yields an honest estimate of the false-positive rate.

\paragraph{Coverage dependency.} Because $\mu$ and $\sigma$ are estimated on
augmented variants, they reflect only the score variability those augmentations
induce. If a real source of variation (cast shadows, light beams, outdoor
background changes through glass) is absent from the augmentation scheme,
$\sigma$ is underestimated and real test scores under that variation exceed
$\theta$ more often than expected. This dependency is the root cause of the
shadow false-positive problem (Section~\ref{sec:p3-results}) and motivates the
shadow and light augmentation introduced below.

\subsubsection{Test phase and anomaly map}

At inference, the query image is processed by the same feature pipeline to
produce its patch vectors. Each patch vector $\phi$ is compared to the memory
bank through nearest-neighbour search; let $m^{\ast}$ be its nearest bank
neighbour, so that the raw patch distance is
$s^{\ast}(\phi) = \lVert \phi - m^{\ast} \rVert_2$. Let
$\mathcal{N}_b(m^{\ast})$ be the set of $b$ nearest neighbours of $m^{\ast}$ within
the bank. Following equation~7 of \cite{Roth2022}, the patch score is re-weighted
as

\begin{equation}
  s(\phi) \;=\; w(\phi)\, s^{\ast}(\phi) , \qquad
  w(\phi) \;=\;
  1 - \frac{\exp \lVert \phi - m^{\ast} \rVert}
  {\sum_{m \in \mathcal{N}_b(m^{\ast})} \exp \lVert \phi - m \rVert} .
  \label{eq:reweight}
\end{equation}

Two points deserve care. First, the distances in the denominator are between the
test patch $\phi$ and the members of $\mathcal{N}_b(m^{\ast})$, \emph{not} between
$m^{\ast}$ and its own neighbours - a detail required for fidelity to the original
formulation. Second, because $m^{\ast}$ is the nearest neighbour, every distance
in the denominator is at least $\lVert \phi - m^{\ast} \rVert$, so the fraction is
at most $1/b$ and the weight is bounded as $w(\phi) \in [\,1 - 1/b,\, 1)$. The
re-weighting therefore never inflates a patch score above its raw distance
$s^{\ast}$: it leaves $s(\phi)$ close to $s^{\ast}$ when $m^{\ast}$ is isolated
(its neighbours are far, the denominator is large, $w \to 1$), and discounts it
by up to a factor $1 - 1/b$ when $m^{\ast}$ sits in a dense region of common
normal patterns (its neighbours are close, $w \to 1 - 1/b$). In relative terms
this promotes patches matched to rare normal prototypes - the configuration most
likely to indicate a genuine anomaly.

The original paper assigns the \emph{image}-level score to the single most
anomalous patch, $\arg\max_{\phi} s^{\ast}(\phi)$, re-weighting only that one
(equations 6-7 of \cite{Roth2022}). The present work instead evaluates
equation~\eqref{eq:reweight} for \emph{every} patch, arranges the scores on the
spatial grid to form an anomaly map, upsamples it bilinearly, and smooths it with a
Gaussian filter; the scalar image score is then the maximum of this smoothed map,
$s_{\max}$, and the binary decision is $s_{\max} > \theta$. This is a deliberate
variant of the paper's single-patch score, adopted for two reasons: the full
per-patch map is required by the spatial gating of
Section~\ref{sec:p3-gating}, which operates on connected components of the map;
and taking the maximum of the smoothed map - rather than of a single raw patch -
is more robust to an isolated noisy patch, since smoothing broadens genuine
anomaly blobs (helping small objects that span few patches) while attenuating
one-pixel spikes. The two formulations coincide when no smoothing is applied and
the maximiser of $w(\phi)\,s^{\ast}(\phi)$ is also the maximiser of $s^{\ast}$.

\subsubsection{Shadow and light augmentation}

The photometric-only memory contains no variants with cast shadows or localised
light excess. This leaves a coverage gap: at test time, a shadow projected on the
floor by a passer-by outside the field of view, or a light beam from a skylight
or a glass door, produces patch features distant from all bank entries and trips
the threshold even though no obstruction is present. The two disturbances are
symmetric - darkening versus brightening - and are addressed symmetrically by
adding synthetic variants of each kind to the augmentation scheme. Three
modalities are modelled for each (elliptical, directional, stripe), covering the
physically plausible patterns in corridor and doorway scenes; all borders are
blurred to avoid unrealistic high-frequency edges.

A disturbance is added to a build variant with probability $p_{\mathrm{sh}}$
(shadow) or $p_{\mathrm{li}}$ (light), and to a calibration variant with
probability $p_{\mathrm{sh}}^{\mathrm{cal}}$ and $p_{\mathrm{li}}^{\mathrm{cal}}$
respectively. A central design choice is to \emph{decouple} the build and
calibration probabilities, because the two phases have opposing requirements:

\begin{itemize}
  \item the \textbf{memory} needs maximum disturbance coverage so that test-time
  shadow/light patches find near neighbours and produce low distances - hence a
  high build probability ($p_{\mathrm{sh}} = p_{\mathrm{li}} = 0.4$);
  \item the \textbf{calibration} needs a stable score distribution to yield a
  reliable $\mu$ and $\sigma$ - hence a low calibration probability
  ($p_{\mathrm{sh}}^{\mathrm{cal}} = p_{\mathrm{li}}^{\mathrm{cal}} = 0.1$).
\end{itemize}

Setting the calibration probability equal to the build probability makes $\sigma$
high-variance: a few unlucky calibration variants with heavy disturbances inflate
$\sigma$, and the $k\sigma$ amplification produces per-camera thresholds that are
either too high (false negatives) or too low (false positives). Decoupling
removes this instability. The two families are independent: a variant may receive
a shadow, a light, both, or neither. The disturbances are deliberately kept out
of the test images used to measure robustness, so that the reported
false-positive rates reflect generalisation to disturbances never seen at build
time.

\subsubsection{Spatial gating via annotated polygons}
\label{sec:p3-gating}

Glass doors introduce a distinct failure mode: the transparent panel makes the
external background visible, and any movement or lighting change beyond the door
generates anomaly blobs in the glass region without any physical obstruction. The
principled remedy exploits a physical-grounding constraint: \emph{any real
obstruction rests on the floor and therefore intersects the door threshold}.

Door-frame polygons (left jamb, right jamb, central mullion, threshold,
architrave) are annotated once on the reference image and rasterised into a
binary \emph{frame mask} $F$. The gating rule operates on the connected components
of the thresholded anomaly map: the map is binarised at $\theta$, its
$8$-connected components are labelled, and a component $C$ is accepted only if

\begin{equation}
  |C \cap F| \;\geq\; \tau ,
  \label{eq:gating}
\end{equation}

where $\tau$ is the minimum overlap, in pixels, between a component and the frame
mask. The final score is the maximum of the continuous anomaly map over the
accepted components, or zero if none survives. A component confined to the glass
panel has zero overlap with $F$ and is rejected, whereas a floor-grounded
obstruction produces a component reaching the threshold or jambs and is accepted.
The same rule applies to opaque doors when the panel exhibits strong illumination
variation: specular highlights stay on the panel, while grounded obstructions
reach the frame.

\paragraph{Structural limitations.} Gating assumes the frame geometry is stable
and visible at test time. An \emph{open door} violates this: the threshold and
jambs change radically, producing an out-of-distribution component inside $F$
(a structural false positive), while an obstacle in the open doorway no longer
intersects the misplaced mask (a structural false negative). A thin-stemmed
grounded object (a floor fan, a sign on a pole) has only a narrow contact with
the frame that may fail the overlap criterion $\tau$. These cases share one root:
reachability of the frame mask by a genuine anomaly component depends on frame
visibility, a property of the scene rather than of the model. They are declared
as limitations of the method; the system is designed and validated for the
nominal closed-door condition, dominant in continuous surveillance.

\subsubsection{Calibration stabilisation: pooled sigma}

The per-camera rule $\theta_i = \mu_i + k\,\sigma_i$ has a structural instability.
The per-camera standard deviation $\sigma_i$ is estimated from a small calibration
sample and carries a relative error on the order of $18\%$; this error propagates
one-to-one into the $k\sigma_i$ term, and across many cameras it produces a
long-tailed distribution of thresholds - a few cameras receive thresholds high
enough to miss all obstructions, others low enough to generate shadow false
positives, even though the underlying scenes are not fundamentally different. The
extreme $\sigma_i$ values are stochastic calibration outliers, not scene outliers:
the same camera re-calibrated with a different seed would receive a normal
threshold.

The fix replaces the per-camera $\sigma_i$ with a single population estimate while
keeping $\mu_i$ local:

\begin{equation}
  \theta_i \;=\; \mu_i + k\,\sigma_{\mathrm{pop}} ,
  \qquad
  \sigma_{\mathrm{pop}} = \operatorname{median}(\sigma_1, \dots, \sigma_N) ,
  \label{eq:pooled}
\end{equation}

where $N$ is the number of cameras. The median is robust to the outlier $\sigma_i$
that caused the tail. Keeping $\mu_i$ per-camera preserves scene-specific
adaptation (the absolute score level of each scene), while pooling $\sigma$
removes the noisy multiplicative term. Under the per-camera rule
$\operatorname{Var}(\theta_i) = \operatorname{Var}(\mu_i) + k^2
\operatorname{Var}(\sigma_i) + 2k\operatorname{Cov}(\mu_i,\sigma_i)$, whereas under
the pooled rule it collapses to $\operatorname{Var}(\theta_i) =
\operatorname{Var}(\mu_i)$, i.e. to the intrinsic scene variability alone.

Pooled sigma is, by construction, a shrinkage (regularisation) estimator: it
trades a small bias - forcing each camera toward the population median - for a
large variance reduction. In the data-starved regime of this thesis (a one-shot
memory from a single image with sparse synthetic calibration) the per-camera
$\sigma_i$ is noise-dominated and the true cross-camera heterogeneity is near
zero, so full pooling is the correct corner. The principled generalisation is
\emph{partial pooling}, $\sigma_i^{\ast} = (1-w_i)\,\sigma_{\mathrm{pop}} +
w_i\,\sigma_i$, with the weight $w_i$ rising as a camera accumulates real normal
frames; full pooling ($w_i = 0$) is the cold-start corner and remains valuable
for newly installed cameras.

\section{Dataset Description}

[Description of the dataset used to build, calibrate, and evaluate the pipeline -
to be completed.]
